% !TeX TS-program = xelatex
% Suiyang Guo — Embedded Systems / Robotics Resume (English)
\documentclass{resume-photo}

\usepackage{xeCJK}
\setCJKmainfont{FandolSong-Regular.otf}[BoldFont=FandolSong-Bold.otf, ItalicFont=FandolKai-Regular.otf]
\usepackage{fontspec}
\usepackage{qrcode}
\usepackage{enumitem}
\usepackage{needspace}
\setlist[itemize]{itemsep=1.5pt, topsep=1.5pt, parsep=1pt}

\RenewDocumentCommand{\ResumeItem}{omoo}{%
  \Needspace{4\baselineskip}%
  \par\vspace{3pt}\noindent
  {\heiti\zihao{-4} #2}%
  \IfValueT{#4}{\hfill{\normalfont\zihao{5} #4}}%
  \par\nopagebreak
  \IfValueT{#3}{{\zihao{5} #3}\par}%
  \vspace{1pt}%
}
\ctexset{section/beforeskip=0.55em, section/afterskip=0.45em}

\ResumeName{Suiyang Guo}
\ResumePhoto{photo.jpg}

\begin{document}

\ResumeContacts{
  +86 158-9334-6463,%
  \ResumeUrl{mailto:1327217178@qq.com}{1327217178@qq.com},%
  \textnormal{Northwestern Polytechnical Univ. | MSc, Control Eng. | 2024.09—2027.06}%
}

\ResumeTitle

\noindent Control Engineering MSc candidate at Northwestern Polytechnical University (Class of 2027), targeting embedded software and robotic systems roles. Proficient in C/C++, STM32, FreeRTOS, embedded Linux and ROS/ROS2, with hands-on experience across firmware, motor control, PCB design, host applications and full-system integration. Built a Yocto/i.MX6ULL robot controller, in-house dual-arm/chassis motor-control systems and real-time VLA deployment pipelines. Contributor to a national key robotics project; holder of 6 patents and 2 software copyrights.

\section{Education}
\ResumeItem{Northwestern Polytechnical University (NPU)}
[\textnormal{MSc, Control Engineering | Robotics \& AI | GPA: 89.56/100 (Top 30\%)}]
[2024.09—2027.06]

\ResumeItem{Zhengzhou University}
[\textnormal{BEng, Automation (Elite Class) | Robotics \& AI | GPA: 3.54/4.0 (Top 20\%)}]
[2020.09—2024.06]

\section{Technical Skills}
\begin{itemize}
  \item \textbf{Programming \& Systems}: C/C++, Python and MATLAB; Linux driver/application development, Git, debugging and multi-module system integration
  \item \textbf{Embedded Linux / RTOS}: Yocto, i.MX6ULL, STM32F4, 8051, Arduino and FreeRTOS; board bring-up, lower-level control, sensor interfaces and motor drivers
  \item \textbf{Robotics}: ROS/ROS2, navigation, SLAM, motion planning, inverse kinematics, point-cloud processing, collision detection and fuzzy-PID control
  \item \textbf{Hardware \& Applications}: Altium Designer/JLCPCB, peripheral circuits, Raspberry Pi, PyQt5 host applications and 3D printing
  \item \textbf{Edge AI \& Performance}: PyTorch/JAX deployment, Triton kernels, CUDA Graphs and RTC; experience integrating AI policies into real-time robot systems
\end{itemize}

\section{Embedded \& Robotics Projects}

\ResumeItem{Series-Parallel Hybrid Wheeled Humanoid Robot}
[\textnormal{National Key Research Project | Key Provincial Lab, Team Member}]
[2024.09—Present]

\begin{itemize}
  \item Integrated a wheeled humanoid comprising an AGV base, parallel mechanism, dual 6-DoF arms and a 2-DoF head
  \item Developed the C++/ROS Noetic master controller for 3D mapping, path planning, object recognition, dynamic obstacle avoidance, point-cloud publishing, collision detection and inverse kinematics
  \item Built a MuJoCo simulation/control-integration environment and developed in-house dual-arm and chassis motor-control systems for actuator and full-robot bring-up
  \item Built a Yocto-based embedded-Linux image for i.MX6ULL, ran roscore and chassis control on-board, and implemented wired-Ethernet ROS multi-machine communication
  \item Developed STM32F4 lower-level position/chassis control and optimized the multi-DoF arm-control scheme, doubling the 6-DoF arm control frequency
\end{itemize}

\ResumeItem{Three Generations of TCM Therapy Devices}
[\textnormal{Challenge Cup Provincial Second Prize | Core R\&D Member}]
[2021.03—2023.09]

\begin{itemize}
  \item Led electronics and software development across three product generations, spanning circuits/PCB, embedded firmware, host applications and system integration
  \item Built STM32/FreeRTOS control systems with fuzzy PID, temperature/image/pressure sensing, Bluetooth communication and multi-task scheduling
  \item Designed STM32 peripheral circuits and PCBs; developed servo/stepper drivers and a smoke-handling module; supported iterative prototype bring-up and debugging
\end{itemize}

\ResumeItem{AI-LLM Voice-Interaction Robot}
[\textnormal{Outstanding Undergraduate Thesis | Solo Developer}]
[2023.12—2024.09]

\begin{itemize}
  \item Independently built a Raspberry Pi 4B/Python/PyQt5/ROS voice robot and host application
  \item Integrated audio capture, ASR, LLM dialogue, speech synthesis, digital-human lip sync and smart-home control; received 2 software copyrights
\end{itemize}

\section{Engineering Experience}

\ResumeItem{IDEA Research: Visincept}
[\textnormal{VLA \& World Model Intern | Robot Systems and Edge AI}]
[2026.04—Present]

\begin{itemize}
  \item Deployed dual-arm VLA policies and migrated the inference service from ROS1 to ROS2, implementing a dual-clock control pipeline with 30 Hz command publishing and 3 Hz model inference
  \item Applied RTC, CUDA Graphs and Triton autotuning to reduce step latency from 256.5 ms to 32.05 ms (8$\times$) and idle end-effector jitter by 81\%
  \item Brought up X-VLA on the in-house dual-arm platform and used firmware EndPoseCtrl IK to reduce the control cycle from 1.4 s to 0.33 s
  \item Built preflight safety checks, continuous-gripper mapping and WebSocket/ROS2 interfaces for model-to-robot integration
\end{itemize}

\ResumeItem{Zhongqin Herui (Shaanxi) Technology Co., Ltd.}
[\textnormal{Technical Partner, Part-time}]
[2025.11—2026.01]

\begin{itemize}
  \item Selected and calibrated sensors and developed 3D mapping/navigation for an autonomous orchard inspection vehicle
  \item Built 3D maps from raw LiDAR and converted them into occupancy grids for navigation
\end{itemize}

\section{Awards \& Intellectual Property}
\begin{itemize}
  \item Challenge Cup Henan Collegiate Academic Sci-Tech Works Competition \hfill Provincial Second Prize (2023)
  \item Collegiate Electrical \& Electronic Engineer Innovation Contest \hfill Central China First Prize (2023)
  \item Siemens Cup Intelligent Manufacturing Challenge \hfill Western Region Second Prize (2022)
  \item Lanqiao Cup, Embedded Design \& Development \hfill Provincial Third Prize (2022)
  \item 4 invention patents, 2 utility-model patents and 2 software copyrights; CET-6: 473
\end{itemize}

\vfill
\begin{center}
  \qrcode[height=2cm]{https://youngyang-gthb.github.io/}\\[4pt]
  {\footnotesize Portfolio: \ResumeUrl{https://youngyang-gthb.github.io/}{youngyang-gthb.github.io}}
\end{center}

\end{document}
